\documentclass{article}
\usepackage{graphicx} % Required for inserting images
\usepackage{amsmath}
\usepackage{algorithmic}
\usepackage{algorithm}
\usepackage[polish]{babel}
\usepackage[utf8]{inputenc}
\usepackage[T1]{fontenc}
\usepackage{float}
\usepackage{listings}

\title{Algorytmy Równoległe - Laboratorium 3}
\author{Adam Staniszewski}

\lstset{frame=tb,
  language=Python,
  aboveskip=3mm,
  belowskip=3mm,
  showstringspaces=false,
  columns=flexible,
  basicstyle={\small\ttfamily},
  numbers=none,
  numberstyle=\tiny\color{gray},
  keywordstyle=\color{blue},
  commentstyle=\color{dkgreen},
  stringstyle=\color{green},
  breaklines=true,
  breakatwhitespace=true,
  tabsize=3
}

\begin{document}

\maketitle
\textbf{Na początku przypomimnijmy co udało się ustalić w ramach laboratorium 1.}

\section*{Wybór problemu}

Wyznaczyć temperaturę w przewodniku elektrycznym, który został nagrzany w wyniku przepływu prądu elektrycznego o stałym natężeniu (powierzchnia boczna jest utrzymywana w stałej temperaturze równej zero).
Zastosować metodę różnicową do równania Poissona:

\[
\frac{\partial^2 T}{\partial x^2} + \frac{\partial^2 T}{\partial y^2} = -\frac{g}{\lambda}
\]

gdzie:
- \( g \) - wydajność źródeł ciepła (stała),
- \( \lambda \) - przewodność cieplna (stała).


\section*{Wyprowadzenie wzoru iteracyjnego}
Warunki brzegowe - powierzchnia boczna jest utrzymywana w stałej temperaturze równej zero.
Równanie Poissona dla rozkładu temperatury zapisujemy jako:
\[
\frac{\partial^2 T}{\partial x^2} + \frac{\partial^2 T}{\partial y^2} = -\frac{g(x, y)}{\lambda}
\]
gdzie:
\begin{itemize}
    \item \( T(x, y) \) - funkcja temperatury,
    \item \( g(x, y) \) - funkcja źródła ciepła,
    \item \( \lambda \) - przewodność cieplna.
\end{itemize}
Przybliżamy drugie pochodne różnicami skończonymi. Dla drugiej pochodnej względem \( x \) mamy:
\[
\frac{\partial^2 T}{\partial x^2} \approx \frac{T_{i+1, j} - 2T_{i, j} + T_{i-1, j}}{dx^2}
\]
Analogicznie, dla drugiej pochodnej względem \( y \):
\[
\frac{\partial^2 T}{\partial y^2} \approx \frac{T_{i, j+1} - 2T_{i, j} + T_{i, j-1}}{dy^2}
\]
Sumujemy teraz oba przybliżenia, co daje:
\[
\frac{T_{i+1, j} - 2T_{i, j} + T_{i-1, j}}{dx^2} + \frac{T_{i, j+1} - 2T_{i, j} + T_{i, j-1}}{dy^2} = -\frac{g_{i, j}}{\lambda}
\]
Zakładając, że siatka jest jednolita, czyli \( dx = dy \), możemy uprościć to wyrażenie do postaci:
\[
\frac{T_{i+1, j} + T_{i-1, j} + T_{i, j+1} + T_{i, j-1} - 4T_{i, j}}{dx^2} = -\frac{g_{i, j}}{\lambda}
\]
Rozwiązując to równanie względem \( T_{i, j} \), otrzymujemy wzór na nową wartość temperatury:
\[
T_{\text{new}}[i][j] = \frac{1}{4} \left( T_{i+1, j} + T_{i-1, j} + T_{i, j+1} + T_{i, j-1} + \frac{g_{i, j} \cdot dx^2}{\lambda} \right)
\]

\section*{Psuedokod algorytmu równoległego z wykorzystaniem metodyki PCAM}
Algorytm równoległy dzieli siatkę na bloki przetwarzane równolegle przez różne procesory z wymianą danych granicznych.

\begin{algorithm}[H]
\caption{Algorytm równoległy (metoda różnic skończonych z MPI)}
\begin{algorithmic}[1]
    \STATE Podziel domenę na $n$ części, gdzie $n$ to liczba procesów.
    \STATE Każdy proces otrzymuje lokalny blok danych, zawierający:
    \STATE \quad Obszar obliczeniowy $T_{\text{local}}[x][y]$ oraz dane potrzebne do obliczeń: $g$, $\lambda$, $dx$, $dy$.
    
    \STATE Uruchom MPI, zdefiniuj $rank$ (numer procesu) oraz $size$ (liczba procesów)
    \STATE Każdy proces otrzymuje swój blok danych, w tym:
    \STATE \quad $g$, $\lambda$, $dx$, $dy$, $T_{\text{local}}[x][y]$
    
    \REPEAT
    \FOR {każdy punkt wewnętrzny $(i, j)$ w lokalnym bloku}
        \STATE Oblicz nową wartość temperatury:
        \[
        T_{\text{new}}[i][j] = \frac{1}{4} (T_{\text{local}}[i+1][j] + T_{\text{local}}[i-1][j] + T_{\text{local}}[i][j+1] + T_{\text{local}}[i][j-1] + \frac{g \times dx^2}{\lambda})
        \]
    \ENDFOR
    
    \STATE Wymień dane graniczne z sąsiadującymi procesami (zbiorczo, w jednym przesłaniu)
    
    \STATE Każdy proces oblicza lokalną różnicę $|T_{\text{new}}[i][j] - T_{\text{local}}[i][j]|$

    \STATE Zbierz różnice na procesie 0 (root) za pomocą $MPI\_Reduce$
    
    \IF {globalna różnica jest mniejsza niż $\epsilon$ (na procesie 0)}
    \STATE Zbieżność osiągnięta, zakończ iterację
    \ENDIF
    \STATE Zaktualizuj macierz $T$: $T_{\text{local}} = T_{\text{new}}$
    \UNTIL{osiągnięcie zbieżności lub maksymalnej liczby iteracji}
    
    \STATE Proces 0 zbiera wyniki z wszystkich procesów za pomocą $MPI\_Gather$

    \STATE Wyświetl wynikową macierz $T$ (na procesie 0)
    
\end{algorithmic}
\end{algorithm}

\textbf{Przejdźmy teraz do treści laboratorium 3}
\section*{Implementacja algorytmu w języku Python}
\begin{lstlisting}
import mpi4py
import numpy as np
import matplotlib.pyplot as plt
from mpi4py import MPI

def main():
    Lx: float = 10.0
    Ly: float = 10.0
    Nx: int = 20
    Ny: int = 20
    g: float = 1.0
    lambda_: float = 1.0

    dx = Lx / (Nx - 1)
    dy = Ly / (Ny - 1)

    comm = MPI.COMM_WORLD
    rank = comm.Get_rank()
    size = comm.Get_size()

    rows_per_process = Ny // size
    extra_rows = Ny % size

    if rank < extra_rows:
        local_rows = rows_per_process + 1
    else:
        local_rows = rows_per_process


    T_local = np.zeros((local_rows, Nx))

    T_local[0, :] = 0
    T_local[-1, :] = 0

    tolerance: float = 1e-6
    error: float = 1.0

    while error > tolerance:
        T_old = T_local.copy()

        for j in range(1, local_rows - 1):
            for i in range(1, Nx - 1):
                T_local[j, i] = (T_old[j + 1, i] + T_old[j - 1, i] + 
                                T_old[j, i + 1] + T_old[j, i - 1] - 
                                (g / lambda_) * (dx ** 2 + dy ** 2)) / 4.0
                
        if rank > 0: 
            comm.Send(T_local[1, :], dest=rank - 1, tag=0)  
            comm.Recv(T_local[0, :], source=rank - 1, tag=1)  

        if rank < size - 1:  
            comm.Send(T_local[-2, :], dest=rank + 1, tag=1)  
            comm.Recv(T_local[-1, :], source=rank + 1, tag=0)

        error = np.max(np.abs(T_local - T_old))
        error = comm.allreduce(error, op=MPI.MAX) 
    T_global = None
    if rank == 0:
        T_global = np.empty((Ny, Nx))

    comm.Gather(T_local, T_global, root=0)

    if rank == 0:
        return T_global
    return None

\end{lstlisting}

\section*{Wyniki}

\begin{figure}[H]
    \centering
    \includegraphics[width=1\linewidth]{obraz.png}
    \caption{Wynikowo heatmapa}
    \label{fig:enter-label}
\end{figure}

\end{document}
